% arara: xelatex
\documentclass[10pt,a4paper]{article}
\usepackage[margin=1.8cm]{geometry}
\usepackage[T1]{fontenc}
\usepackage[utf8]{inputenc}
\usepackage[portuguese]{babel}
\usepackage{newcomputermodern}
\usepackage{microtype}
\usepackage{xcolor}
\usepackage[hidelinks]{hyperref}
\usepackage{enumitem}
\usepackage{titlesec}
\usepackage{tabularx}
\usepackage{array}
\usepackage{setspace}
\usepackage{ragged2e}
\usepackage{parskip}
\usepackage{fancyhdr}

\pagestyle{fancy}
\fancyhf{}
\renewcommand{\headrulewidth}{0pt}
\renewcommand{\footrulewidth}{0pt}
\setlength{\headheight}{14pt}
\pagenumbering{gobble}

\definecolor{accent}{HTML}{1F4E79}
\definecolor{textgray}{HTML}{404040}

\hypersetup{
  colorlinks=true,
  urlcolor=accent,
  linkcolor=accent,
  citecolor=accent,
  pdftitle={Heliton Martins Reis Filho - CV},
  pdfauthor={Heliton Martins Reis Filho},
  pdfsubject={Curriculum Vitae},
  pdfkeywords={CV, Currículo, Engenheiro de Software, Machine Learning, IA, LLM, RAG, QSAR},
  pdfcreator={LaTeX},
}

\titleformat{\section}{\Large\bfseries\color{accent}}{}{0pt}{}[\vspace{0.22em}\titlerule\vspace{0.22em}]
\titlespacing*{\section}{0pt}{1.1em}{0.42em}
\setlength{\parindent}{0pt}
\setlength{\parskip}{0.15em}
\setlist[itemize]{leftmargin=1.2em,itemsep=0.16em,topsep=0.14em,parsep=0pt,partopsep=0pt}

\newcommand{\entry}[4]{%
  \textbf{\sffamily\large #1} \hfill #2 \\
  {\itshape #3} \hfill {\small #4} \\[-0.75em]
}

\newcommand{\subentry}[2]{%
  \hspace{0.4em}\textbf{\sffamily #1} \hfill {\small #2}\\[-0.55em]
}

\newcommand*{\itemdesc}[2]{\item \textbf{\sffamily #1:}~{\small #2}}

\newcommand*{\link}[2]{\href{#1}{{\texttt{#2}}}}

\begin{document}

{\Huge \textbf{Heliton Martins Reis Filho}}\\[0.3em]

{\large Engenheiro de \textit{Software} e IA/ML | Computação Científica \& Inovação Estratégica}\\[0.4em]

{\ttfamily
\begin{tabularx}{\textwidth}{@{}X r@{}}
UFMG | CEFET-MG & Belo Horizonte, MG, Brasil \\
\href{https://github.com/hellmrf}{github.com/hellmrf} & \href{https://hellmrf.dev.br}{hellmrf.dev.br} \\
\href{https://www.linkedin.com/in/helitonmrf/}{linkedin.com/in/helitonmrf} & \href{https://wa.me/5533999429795}{(33) 99942-9795} \\
\href{https://lattes.cnpq.br/2197799814329542}{lattes.cnpq.br/2197799814329542}& \href{mailto:helitonmrf@gmail.com}{helitonmrf@gmail.com} \\
\end{tabularx}
}

\section*{Resumo Profissional}
Engenheiro de \textit{Software} e Empreendedor de Base Tecnológica com um perfil multidisciplinar raro, unindo o pragmatismo de mais de uma década em engenharia de sistemas ao profundo rigor matemático da Química Computacional. Apaixonado por aprendizado contínuo, atuo na interseção entre ciência profunda, arquitetura de \textit{software} e estratégia de negócios, entregando valor através de três pilares:

\begin{itemize}[leftmargin=1.5em, itemsep=0.4em]
    \item \textbf{Engenharia de \textit{Software} Versátil:}
    Ampla fluência tecnológica (10+ linguagens, incluindo Rust, Python e TypeScript) e domínio de múltiplas eras da infraestrutura (de servidores Linux e VPS \textit{bare-metal} a ecossistemas \textit{Cloud-Native}, Docker, CI/CD e \textit{Serverless}).

    \item \textbf{Ciência de Dados, IA e Rigor Matemático:}
    Sólida base analítica aplicada a \textit{hard sciences}. Extensa experiência em software científico (QSAR), redução de dimensionalidade, validação estatística e aplicação de Redes Neurais para elucidação de sistemas complexos.

    \item \textbf{Visão Estratégica e Inovação (\textit{Lean}):}
    Liderança técnica com mentalidade de produto. Especialista em aplicar ciclos ágeis (construir-medir-aprender) e alinhar o desenvolvimento tecnológico a análises de viabilidade econômica, macropolítica e necessidades reais do mercado para construir \textit{startups} e soluções de alto impacto.
\end{itemize}


\section*{Competências Técnicas}
\begin{itemize}
\itemdesc{IA / Machine Learning}{LLMs, RAG, busca semântica, \textit{embeddings} (\textit{Transformers}), \textit{fine-tuning}, modelos locais, bancos vetoriais, redes neurais, regressão e classificação, validação estatística, aprendizagem (não-)supervisionada}
\itemdesc{Linguagens}{Python, Rust, TypeScript/JavaScript, SQL, Julia, C++, Fortran, MATLAB, \LaTeX, PHP, HTML/CSS}
\itemdesc{Backend e Web}{Next.js, NestJS, Node.js, FastAPI, Flask, APIs REST, PostgreSQL, bancos NoSQL e vetoriais (Pinecone)}
\itemdesc{Infra e DX}{Git/GitHub, Docker, CI/CD, Linux, AWS/GCP/Azure, MLOps, \textit{serverless}}
\end{itemize}

\section*{Experiência Profissional}

\entry{Coral Digital Soluções Inteligentes}{Mai. 2024 -- Atual}{Fundador, CEO \& CTO}{\link{https://coraldigital.com.br}{coraldigital.com.br}}
\begin{itemize}\footnotesize
  \item Fundei e lidero, estratégica e tecnicamente, a Coral Digital, atuando em duas frentes: produtos próprios de IA aplicada e consultoria especializada em ML/IA para empresas;
  \item Aplico metodologia enxuta (ciclos construir--medir--aprender, com contato direto com clientes desde o MVP) no desenvolvimento de novos produtos;
  \item Prestei serviços sob demanda, incluindo um classificador semântico de licitações (\textit{embeddings} + ML clássico) e \textit{chatbots} conversacionais com LLMs para clientes diversos.
\end{itemize}

\subentry{Licita.AI --- SaaS de Inteligência em Licitações Públicas}{Set. 2025 -- Atual \textbar{} \link{https://licitaai.coraldigital.com.br}{licitaai.coraldigital.com.br}}
\begin{itemize}[leftmargin=2.2em]\footnotesize
  \item Idealizo e lidero, como CEO \& CTO, o Licita.AI, incubado no Berçário de Startups do DQ/UFMG, com mentoria estruturada e parceiros estratégicos desde o MVP;
  \item Arquitetei o produto em 4 microsserviços: aplicação web \textit{fullstack} em Next.js, dois \textit{workers multithreaded} em Rust e um em Python para a \textit{pipeline} de dados (\textit{serverless});
  \item Implementei busca semântica de licitações com \textit{embeddings} e banco vetorial, e pré-análise de viabilidade de editais via LLM, com extração estruturada de dados a partir dos anexos;
  \item Conduzi a migração de arquitetura de dados de PostgreSQL + Pinecone para PostgreSQL único, reduzindo superfície de infraestrutura e custo operacional.
\end{itemize}

\subentry{Responde.Chat --- SaaS de Automação de Atendimento via WhatsApp}{2024 -- 2025}
\begin{itemize}[leftmargin=2.2em]\footnotesize
  \item Projetei e desenvolvi a solução completa: API própria de WhatsApp \textit{self-hosted} (Evolution API), \textit{backend} em NestJS, \textit{frontend} em Next.js e funções/\textit{webhooks} em Python/FastAPI;
  \item Entreguei um produto configurável pelo próprio WhatsApp, eliminando a necessidade de painel administrativo separado para o cliente final.
\end{itemize}

\vspace{0.7em}

\entry{Companhia Brasileira de Lítio (CBL)}{Mar. 2025 -- Abr. 2025}{Estagiário em Pesquisa \& Desenvolvimento}{Divisa Alegre, MG}
\begin{itemize}\footnotesize
  \item Estágio obrigatório no setor de P\&D de uma planta industrial de beneficiamento de Lítio, na interface entre química, processos industriais e pesquisa aplicada.
\end{itemize}

\vspace{0.7em}

\entry{Privacy Point}{Ago. 2024 -- Jan. 2025}{Desenvolvedor Full-Stack}{\link{https://privacypoint.com.br}{privacypoint.com.br}}
\begin{itemize}\footnotesize
  \item Único desenvolvedor \textit{full-stack} do produto (ao lado do CTO), entregando funcionalidades voltadas à conformidade com a LGPD em plataforma PHP/Laravel sobre VPS;
  \item Desenvolvi e implantei \textit{chatbot} conversacional especializado em privacidade e proteção de dados (LLM + RAG + banco vetorial).
\end{itemize}

\section*{Experiência em Pesquisa Científica}

\entry{LACC/UFMG --- Laboratório de Química Computacional Aplicada}{2019 -- Atual (8 anos)}{Pesquisador Voluntário e Líder de Desenvolvimento de Software}{DQ/UFMG}
\begin{itemize}\footnotesize
  \item Liderei o desenvolvimento do QSARModelingPy: cálculo de descritores moleculares, tratamento de dados, filtragem e seleção de variáveis (variância, correlação, autocorrelação, Kennard-Stone, OPS, algoritmo genético) e validação estatística (interna, cruzada e externa);
  \item Reimplementei módulos críticos em Rust e Julia para comparação de performance com a versão original em Python;
  \item Concebi e implementei um método inédito baseado em geometria computacional para cálculo de descritores moleculares, tema do meu Trabalho de Conclusão de Curso;
  \item Publico e mantenho ferramentas de código aberto (HullQSAR, qsarmodeling, QSARModelingPy, maligner) que conectam pesquisa científica a \textit{software} reutilizável pela comunidade;
  \item Atuo em modelagem molecular, cálculos quânticos em \textit{clusters} (ORCA, GAMESS, Psi4, OpenBabel, RDKit) e arquitetura/empacotamento de aplicações científicas em ambiente Linux.
\end{itemize}

\vspace{0.7em}

\entry{PICQ/UFMG --- Laboratório de Problemas Inversos e Cinética Química}{2020 -- 2024 (4 anos)}{Bolsista de Iniciação Científica (CNPq)}{DQ/UFMG}
\begin{itemize}\footnotesize
  \item Apliquei redes neurais artificiais e métodos isoconversionais/isotérmicos (Arrhenius, Vyazovkin) a dados de análise térmica (DTG/DTA) para elucidação de mecanismos de decomposição;
  \item Reimplementei as metodologias legadas do grupo (originalmente em MATLAB) em Julia, Fortran e Python, com foco em \textit{performance};
  \item Fui responsável metodológico e de desenvolvimento de \textit{software} em projetos interdisciplinares do PPG em Inovação Tecnológica, incluindo: estudo de espectroscopia Raman + redes neurais para detecção de COVID-19 com \textbf{97\% de acurácia}; otimização via RNA do tempo de duração de processos na Justiça do Trabalho; e pré-análise de processos judiciais sobre proteção de dados com LLMs e ML clássico.
\end{itemize}

\section*{Formação Acadêmica}
\entry{Universidade Federal de Minas Gerais (UFMG)}{2019 -- jul. 2026}{Bacharelado em Química com ênfase em Química Computacional}{Belo Horizonte, MG}
\begin{itemize}\footnotesize
  \item Disciplinas complementares:
  \begin{itemize}
      \item Programação e Desenvolvimento de Software; Introdução à Lógica Computacional; Álgebra Linear Computacional; Matemática Discreta;
      \item Introdução à Quimiometria; Introdução à Ciência de Dados; Métodos Computacionais e Estatísticos em Química; Métodos da Física Teórica;
      \item Criação de Empresas de Base Tecnológica; Vivência em \textit{Startups} de Base Tecnológica; Empreendimentos em Informática.
    \end{itemize}
\end{itemize}

\vspace{0.5em}

\entry{Centro Federal de Educação Tecnológica de Minas Gerais (CEFET-MG)}{2016 -- 2018}{Técnico em Química}{Timóteo, MG}
\begin{itemize}\footnotesize
  \item Estágio no Núcleo de Estudos em Física e Química Computacional (NEFICQ), na linha de planejamento de fármacos e \textit{docking} molecular.
\end{itemize}

\section*{Projetos de Código Aberto}
\begin{itemize}
  \item \link{https://github.com/lacc-ufmg/qsarmodeling}{QSAR Modeling} --- aplicação em Rust para modelagem e validação para QSAR/QSPR.
  \item \link{https://github.com/hellmrf/HullQSAR}{HullQSAR} --- biblioteca para cálculo geométrico de descritores moleculares.
  \item \link{https://github.com/lacc-ufmg/QSARModelingPy}{QSARModelingPy} --- biblioteca em Python para construção e validação de modelos QSAR/QSPR.
  \item \link{https://github.com/hellmrf/maligner}{maligner} --- ferramenta de alinhamento molecular.
\end{itemize}

\section*{Divulgação Científica e Ensino}
\begin{itemize}
  \item Produtor (Coordenador de Transmissão) do \href{https://repositorio.ufmg.br/bitstream/1843/50717/1/I%20Congresso%20Nacional%20de%20Inova%C3%A7%C3%A3o%20e%20Populariza%C3%A7%C3%A3o%20da%20Ci%C3%AAncia%20a%C3%A7%C3%B5es%20durante%20a%20covid-19.pdf}{I Congresso Nacional de Inovação e Popularização da Ciência} (2020);
  \item Fundador do canal \link{https://www.youtube.com/channel/UCsXyoSkpyZB741fkBKeEovQ}{Programação Popular}, com mais de 1 milhão de visualizações desde 2013, focado no ensino acessível de programação;
  \item Produtor do Curso de Lógica de Programação e Algoritmos em Portugol Studio, adotado como material de apoio em Institutos Federais e Escolas Técnicas brasileiras.
\end{itemize}

\section*{Reconhecimentos}
\begin{itemize}
  \item \textbf{Menção Honrosa} na \href{https://obesq.obquimica.org}{Olimpíada Brasileira do Ensino Superior de Química}.
  \item \textbf{Medalha de Ouro} na \href{https://www.meta2017.cefetmg.br}{XXVII Mostra Específica de Trabalhos e Aplicações} do CEFET-MG.
\end{itemize}

\section*{Liderança e Atuação Institucional}
\begin{itemize}
  \item Representante discente em órgãos colegiados da UFMG (Conselho Universitário, CEPE, Congregação do ICEx, Conselho Fiscal da FUMP, Câmara Departamental de Química);
  \item Atuação executiva em entidades estudantis: DCE-UFMG (Diretor Administrativo-Financeiro), DAICEx (Secretário Geral) e CEQ (Presidente).
\end{itemize}

\section*{Idiomas}
Português (nativo) \textbar{} Inglês (fluente) \textbar{} Italiano (básico)

\begin{flushright}\scriptsize%
Atualizado em \today.%
\end{flushright}

\end{document}
